\pdfmapfile{+symbols.map}
\pdfmapfile{+cm.map}
\pdfmapfile{+cmextra.map}
\documentclass[11pt]{article}
\usepackage[margin=1in]{geometry}
\usepackage{amsmath,amssymb,amsthm}
\usepackage[hidelinks]{hyperref}
\newtheorem{theorem}{Theorem}
\newtheorem{lemma}{Lemma}
\newcommand{\N}{\mathbb N}
\title{Positive Tutte Evaluations and an Explicit Finite Count}
\author{Galaxy-0}
\date{October 5, 2026}
\begin{document}
\maketitle
\begin{abstract}
For every finite undirected multigraph $G$ and integer $q\ge2$, we prove
$T_G(q+1,q-1)\ge q^{r(E)}>0$. We construct a finite family of decorated
spanning subgraphs whose cardinality is this evaluation and give a terminating
algorithm to enumerate it. This proves the positivity and effective-count
assertions in conjecture 00000000585, with no prime-power restriction needed.
\end{abstract}

\section{Statement and conventions}
The published conjecture states: ``For every prime power $q$,
$T_G(q+1,q-1)>0$ and yields an effective count (an arithmetization of the
combinatorial positivity interpretation).'' Its preceding descriptive sentence
mentions log-concavity but specifies no sequence, inequality, or additional
log-concavity assertion. The theorem below addresses exactly the stated
positivity and counting assertions; it does not assert an unspecified
log-concavity property.

Let $G=(V,E)$ be a finite undirected graph. Loops, parallel edges, isolated
vertices, disconnected graphs, and the graph with no vertices are permitted.
Edges are individually labeled; consequently parallel edges determine different
subsets of $E$. For $A\subseteq E$, write $c(A)$ for the number of connected
components of the spanning graph $(V,A)$, including isolated vertices, and put
\[
 r(A)=|V|-c(A),\qquad \nu(A)=|A|-r(A).
\]
For the empty vertex set the component count is zero. In particular
$r(\varnothing)=0$. We use the standard rank-subset definition of the Tutte
polynomial:
\begin{equation}\label{eq:tutte}
 T_G(x,y)=\sum_{A\subseteq E}
 (x-1)^{r(E)-r(A)}(y-1)^{|A|-r(A)}.
\end{equation}
All exponents in this polynomial are nonnegative integers. Indeed, adding
edges cannot increase the number of components, so $r(A)\le r(E)$. Starting
with $|V|$ isolated vertices, each inserted edge decreases the component count
by at most one, so $r(A)\le|A|$. These observations cover loops and parallel
edges without any change. As usual a zeroth power is one, including $0^0=1$.

\newpage
\section{Positivity}
\begin{theorem}\label{thm:positive}
For every finite undirected multigraph $G$ and every integer $q\ge2$,
\[
 T_G(q+1,q-1)=\sum_{A\subseteq E}
 q^{r(E)-r(A)}(q-2)^{\nu(A)}\ge q^{r(E)}>0.
\]
\end{theorem}
\begin{proof}
Substitute $x=q+1$ and $y=q-1$ into \eqref{eq:tutte}. Since $q\ge2$, the
bases $q$ and $q-2$ are nonnegative. Every summand is therefore nonnegative.
For $A=\varnothing$, the summand is
$q^{r(E)-0}(q-2)^0=q^{r(E)}$, which is strictly positive. Keeping that
summand and discarding the other nonnegative summands proves the assertion.
\end{proof}
If $q=p^k$ is a prime power, with $p$ a prime and $k\ge1$, then $q\ge2$.
The conjectured positivity is thus an immediate instance of
Theorem~\ref{thm:positive}. The stronger conclusion also includes integers
that are not prime powers.

\section{An explicit counting interpretation}
For an integer $s\ge0$, let $[s]=\{0,\ldots,s-1\}$. The set of words of
length $d$ over this alphabet is $[s]^d$, the set of functions from $[d]$ to
$[s]$. Its cardinality is $s^d$. In particular there is one empty word over
an empty alphabet and no positive-length word over that alphabet.

Define the finite set
\begin{equation}\label{eq:objects}
 \mathcal C(G,q)=\coprod_{A\subseteq E}
 \bigl(\{A\}\times[q]^{r(E)-r(A)}\times[q-2]^{\nu(A)}\bigr).
\end{equation}
An object is a selected spanning subgraph, a $q$-color word in its
rank-deficiency slots, and a $(q-2)$-color word in its cycle-nullity slots.
The selected edge set is part of the object, so contributions from different
subsets cannot be confused. No equivalence of colorings is taken.

\begin{theorem}\label{thm:count}
For every $G$ and every integer $q\ge2$,
\[
 |\mathcal C(G,q)|=T_G(q+1,q-1).
\]
The family $\mathcal C(G,q)$ is nonempty and has a finite, explicit
enumeration algorithm.
\end{theorem}
\begin{proof}
The summand indexed by $A$ in \eqref{eq:objects} has cardinality
$q^{r(E)-r(A)}(q-2)^{\nu(A)}$, by the product rule for cardinalities.
Taking the disjoint union gives the sum in Theorem~\ref{thm:positive}.
The empty edge set supplies $q^{r(E)}$ objects, so the family is nonempty.

To enumerate it, order the vertices and individually labeled edges. Compute
$c(E)$ by graph search. Enumerate the $2^{|E|}$ Boolean edge-selection masks.
For each mask, compute $c(A)$ by graph search, calculate the two displayed
nonnegative lengths, and enumerate the words over the corresponding finite
alphabets in lexicographic order. Output the mask and the two words. Every
loop of this procedure has a specified finite bound. Each valid object
appears exactly once, because its mask and each word have unique encodings.
Thus this is a terminating effective enumeration. Counting instead of
materializing the output gives the same evaluation. No polynomial-time
claim is made.
\end{proof}
At $q=2$, precisely the acyclic subsets contribute: if $\nu(A)>0$, their
contribution is zero; if $\nu(A)=0$, it is $2^{r(E)-r(A)}$. The empty
subset still contributes positively. This also explains why the endpoint
$q=2$ does not present an exception.

\section{Lean definitions and scope}
The accompanying \texttt{Tutte585.lean} uses Lean 4.31.0 and its standard
library only. The complete Lake project and executable examples compile
successfully. The main theorem's audited dependencies are \texttt{propext},
\texttt{Classical.choice}, and \texttt{Quot.sound}; no additional axiom
or proof placeholder occurs.
A graph has a finite vertex type $\operatorname{Fin}(n)$ and a list of pairs
of vertices. List positions are edge labels, so repeated endpoint pairs
represent parallel edges. Equal endpoints represent loops.

The inductive relation \texttt{Connected} is generated by reflexivity,
edges, symmetry, and transitivity. The definition \texttt{componentCount}
counts the least vertex in each connected component. Each nonempty component
has exactly one such least vertex, so this is precisely $c(A)$.
The definition \texttt{rank} is $n-\texttt{componentCount}$, and
\texttt{tutte} is exactly the rank-subset expansion above, evaluated in the
natural numbers for $x,y\ge1$. Only $x=q+1$, $y=q-1$ with $q\ge2$ is claimed
to represent the usual polynomial evaluation. Natural subtraction in the
exponents equals integer subtraction by the graph inequalities proved in
Section~1; these basic rank inequalities are not separately formalized in
this file.

The formal proof establishes empty-graph connectivity as equality and
$r(\varnothing)=0$, proves the lower bound and strict positivity, and proves
that the explicitly defined decorated-subgraph enumeration has the same
length as the evaluation and has no duplicate objects. Decorations are represented by pairs of integers
less than the corresponding powers; these are the ordinary finite word
codes. Boolean masks preserve edge identity throughout enumeration.
The main theorem is \texttt{conjecture\_00000000585}.

The formal component count is executable. A finite cut test enumerates
all vertex subsets and asks whether every edge-closed subset containing $u$
also contains $v$. The file proves that this test is equivalent to
\texttt{Connected}: connectivity preserves membership in a closed cut;
conversely, the component of $u$ itself is an edge-closed cut containing $u$.
Classical reasoning is used only inside this correctness proof. The decision
procedure, component count, rank, and object enumeration are all computable
Lean definitions. This exhaustive cut procedure is deliberately simple and
may be exponential in the number of vertices; graph search gives the more
practical implementation described above.

There are no added axioms or proof placeholders. The final source prints
its theorem's logical dependencies for audit. In particular, ordinary
classical logical dependencies are not graph-theoretic assumptions.

\section*{Source}
The Last Math Competition, conjecture 00000000585, revision
\texttt{d7d4bc9f6918}. \href{https://github.com/The-Last-Math-Competition/The-Last-Math-Competition/blob/d7d4bc9f6918401d3fb22c6e473df63eba5dbb78/conjectures/00000000585.md}{Published statement at the exact pinned revision}.
\end{document}
